% !TeX root = ../nc-quadric-classification.tex
% 2026-11-21 / v1.34
\section{次数別完備化による Noetherian 性の貼合せ}
\label{sec:formal-noetherian-patching}

\paragraph{主張の範囲.}
以下は可換 Noetherian 基底上の非可換次数付き代数と
正の向きの $\mathbb Z$ 代数に対する、証明済みの環論的補題である。
閉部分上とその開補集合上で実際の代数族が Noetherian であることを仮定する。
幾何学的ファイバーの Noetherian 性だけを仮定する主張ではない。
各次数成分の基底上の有限性を用いるが、平坦性、射影性、二次表示、
成長次数、導来同値は仮定しない。新規性は未確認である。

\subsection{中央変数と次数付き有限生成}

\begin{lemma}[中央変数に対する Hilbert 基底論法]
\label{lem:patch-central-variable}
$A$ を環、$M$ を右 Noetherian $A$ 加群とする。
$t$ を中央変数とすれば $M[t]$ は右 Noetherian $A[t]$ 加群である。
内部次数を持つ $A,M$ について、Noetherian 性を次数付き部分加群だけに対して
仮定・結論する形も成立する。この場合 $t$ の内部次数を $0$ とする。
同じ結論は、冪等元 $e_i$ を持つ $\mathbb Z$ 代数 $C$ と
右 $C$ 加群の圏にも成立する。この場合、$t$ は頂点を変えず、
$C[t]_{ij}=C_{ij}[t]$ と定める。
\end{lemma}

\begin{proof}
$N\subset M[t]$ を部分加群とする。$n\geq0$ に対して
$L_n\subset M$ を、$N$ に属し $t$ 次数が $n$ 以下である多項式の
$t^n$ の係数全体とする。これは $A$ 部分加群であり、
$t$ 倍により $L_n\subset L_{n+1}$ である。
$M$ が Noetherian なので、この列はある $n_0$ から安定する。
$0\leq n\leq n_0$ に対して $L_n$ の有限生成元を取り、
それぞれを $N$ の $t$ 次数 $n$ 以下の多項式へ持ち上げる。
この有限個の多項式は $N$ を生成する。
実際、$v\in N$ の最高次係数を消すため、$\deg_t v\leq n_0$ のときは
その次数の持上げを用い、$\deg_t v>n_0$ のときは $n_0$ 次の持上げの
適当な $t$ 冪倍を用いる。差の次数が下がるので、次数に関する帰納法が使える。

内部次数を持つ場合、次数付き $N$ から作る $L_n$ は次数付き部分加群なので、
同じ証明が使える。
$\mathbb Z$ 代数の場合も $L_n$ は $M$ の右部分加群である。
有限生成元を各頂点に同次に取り、多項式の持上げも同じ頂点内で選べば、
まったく同じ帰納法が成り立つ。変数を一つずつ付加することで、
有限個の中央変数についても結論を得る。
\end{proof}

\begin{lemma}[次数付き Noetherian 性から通常の Noetherian 性へ]
\label{lem:patch-graded-to-ungraded}
$A=\bigoplus_{d\geq0}A_d$ を非負次数付き代数とし、
$M$ を下に有界な次数付き右 $A$ 加群とする。
$M$ の次数付き部分加群がすべて有限生成なら、$M$ は
次数を忘れても Noetherian である。
\end{lemma}

\begin{proof}
$N\subset M$ を次数付きとは限らない部分加群とする。
内部次数による増大フィルトレーションを入れると、
$\operatorname{gr} N\subset\operatorname{gr} M=M$ は次数付き部分加群である。
その有限個の同次生成元を $u_1,\ldots,u_s\in N$ に持ち上げる。
任意の $x\in N$ の最高内部次数項を $u_\ell$ の右倍で消せる。
差の最高内部次数は真に下がる。$M$ が下に有界であるから
この操作は有限回で終了し、$u_1,\ldots,u_s$ が $N$ を生成する。
\end{proof}

\subsection{次数ごとの完備性による持上げ}

\begin{lemma}[完備基底上の持上げ]
\label{lem:patch-complete-lifting}
$R$ を可換 Noetherian 環、$I=(f_1,\ldots,f_r)\subset R$ をイデアルとし、
$R$ は $I$ 進完備かつ分離的であると仮定する。
$A=\bigoplus_{d\geq0}A_d$ を中央 $R$ 代数とし、各 $A_d$ は有限 $R$ 加群とする。
このとき $A/IA$ が右 Noetherian なら $A$ も右 Noetherian である。

より一般に、$M$ を各成分が有限 $R$ 加群である次数付き右 $A$ 加群とする。
$M/IM$ が次数付き右 $A/IA$ 加群として Noetherian なら、
$M$ も次数付き右 $A$ 加群として Noetherian である。
$M$ が下に有界なら、次数を忘れても Noetherian である。
\end{lemma}

\begin{proof}
各有限 $R$ 加群 $A_d,M_d$ は $I$ 進完備かつ分離的である。
これは有限加群の完備化が $-\otimes_R\widehat R_I$ と一致することと
$R=\widehat R_I$ から従う
\cite[Tag 00MA]{StacksCompletionPatching}。

$\overline A=A/IA$、$\overline M=M/IM$ と置く。
内部次数と $I$ 進次数を持つ対象に対し、自然な全射
\begin{equation}\label{eq:patch-gr-surjections}
 \overline A[t_1,\ldots,t_r]\twoheadrightarrow\operatorname{gr}_I A,
 \qquad
 \overline M[t_1,\ldots,t_r]\twoheadrightarrow\operatorname{gr}_I M
\end{equation}
がある。$t_\alpha$ は内部次数 $0$、$I$ 進次数 $1$ であり、
$f_\alpha$ の初項へ写す。これらが全射である理由は、
$I^p A_d$ と $I^p M_d$ が、次数 $p$ の $f_\alpha$ の単項式を
係数とする和で生成されることである。
ここでは全射だけを用い、平坦性を仮定しないので同型とは主張しない。
補題~\ref{lem:patch-central-variable} より、
$\operatorname{gr}_I M$ は右 $\operatorname{gr}_I A$ 加群として Noetherian である。

$N\subset M$ を次数付き部分加群とする。
$F^pN=N\cap I^pM$ による誘導フィルトレーションを入れれば
$\operatorname{gr}_F N\subset\operatorname{gr}_I M$ は二重次数付き部分加群である。
有限個の二重同次生成元を取り、それぞれを
\[
 u_\ell\in N_{d_\ell}\cap I^{a_\ell}M_{d_\ell}
 \qquad (1\leq\ell\leq s)
\]
に持ち上げる。これらが $N$ を生成することを示す。

$x\in N_d$ を固定する。$I$ 進次数 $p=0,1,\ldots$ ごとに初項を消せるので、
$a_\ell\leq p$ のとき
\[
 c_{\ell,p}\in I^{p-a_\ell}A_{d-d_\ell},
\]
$a_\ell>p$ のとき $c_{\ell,p}=0$ を選んで、
\begin{equation}\label{eq:patch-approximation}
 x-\sum_{p=0}^{q}\sum_{\ell=1}^s u_\ell c_{\ell,p}
 \in N_d\cap I^{q+1}M_d
 \qquad(q\geq0)
\end{equation}
とできる。存在しない負次数成分は $0$ とする。
固定した $\ell$ について $p-a_\ell\to\infty$ であるから、
$c_\ell=\sum_{p\geq0}c_{\ell,p}$ は有限 $R$ 加群
$A_{d-d_\ell}$ 内で収束する。$M_d$ の分離性と
積の $I$ 進連続性から、\eqref{eq:patch-approximation} の極限は
\[
 x=\sum_{\ell=1}^s u_\ell c_\ell
\]
である。したがって $N$ は有限生成である。
この証明は $N$ 全体の完備性を仮定せず、
固定した内部次数の係数加群の完備性だけを使っている。

$M=A$ とすれば、$A$ のすべての次数付き右イデアルが有限生成である。
補題~\ref{lem:patch-graded-to-ungraded} により通常の右 Noetherian 性を得る。
一般の下に有界な $M$ も同様である。
\end{proof}

\begin{remark}[全体の完備化との区別]
一般の $R$ に対して以下で使う代数は
\[
 A^{\mathrm{deg}\wedge}
 :=A\otimes_R\widehat R_I
 =\bigoplus_{d\geq0}(A_d\otimes_R\widehat R_I)
\]
である。これは各内部次数を個別に完備化した次数付き代数であり、
無限個の内部次数にまたがる元を含み得る環
$\varprojlim_n A/I^nA$ と同一視しない。
補題~\ref{lem:patch-complete-lifting} の収束は一つの内部次数内で行われるため、
$A^{\mathrm{deg}\wedge}$ 全体が $I$ 進完備である必要はない。
\end{remark}

\subsection{閉部分と開補集合からの貼合せ}

\begin{theorem}[Noetherian 性の形式的貼合せ]
\label{thm:patch-noetherian}
$R$ を可換 Noetherian 環、$I=(f_1,\ldots,f_r)\subset R$ とする。
$A=\bigoplus_{d\geq0}A_d$ を各 $A_d$ が有限 $R$ 加群である中央 $R$ 代数とする。
次を仮定する。
\begin{enumerate}
 \item $A/IA$ は右 Noetherian である。
 \item 各 $\alpha$ について $A\otimes_R R_{f_\alpha}$ は右 Noetherian である。
\end{enumerate}
このとき $A$ は右 Noetherian である。左についても同じ結論が成り立つ。

より一般に、各成分が有限 $R$ 加群である次数付き右 $A$ 加群 $M$ について、
$M/IM$ とすべての $M_{f_\alpha}$ がそれぞれ次数付き Noetherian なら、
$M$ も次数付き Noetherian である。
$M$ が下に有界なら通常の右加群としても Noetherian である。
\end{theorem}

\begin{proof}
$\widehat R=\widehat R_I$ と置く。
$\widehat R$ は Noetherian かつ $I\widehat R$ 進完備であり、
$\widehat R/I\widehat R=R/I$ である
\cite[Tags 031C, 05GH]{StacksCompletionPatching}。
$A^{\mathrm{deg}\wedge}/IA^{\mathrm{deg}\wedge}=A/IA$ なので、
補題~\ref{lem:patch-complete-lifting} により
$A^{\mathrm{deg}\wedge}$ は右 Noetherian である。

有限直積
\begin{equation}\label{eq:patch-faithful-cover}
 R^{\ast}=\widehat R_I\times\prod_{\alpha=1}^r R_{f_\alpha}
\end{equation}
を取る。完備化の平坦性 \cite[Tag 00MB]{StacksCompletionPatching} と局所化の平坦性により
$R\to R^{\ast}$ は平坦である。
$\operatorname{Spec}\widehat R_I\to\operatorname{Spec} R$ の像は、
$\widehat R_I/I\widehat R_I=R/I$ より $V(I)$ を含む。
残りは $D(f_\alpha)$ で覆われるから、スペクトル上で全射である。
したがって $R\to R^{\ast}$ は忠実平坦である
\cite[Tag 00HQ]{StacksCompletionPatching}。
なお $I=R$ の場合は開被覆だけを用い、$I=0$ の場合は仮定 (1) で結論が出る。

$A\otimes_R R^{\ast}$ は右 Noetherian 環の有限直積である。
$J\subset A$ を任意の右イデアルとすると、平坦性により
$J\otimes_R R^{\ast}$ は $A\otimes_R R^{\ast}$ の右イデアルであり、有限生成である。
その有限個の生成元をテンソルの有限和として書き、
現れる有限個の $J$ の元で生成される右イデアルを $J_0$ とする。
すると
\[
 (J/J_0)\otimes_R R^{\ast}=0.
\]
忠実平坦性より $J=J_0$ である。したがって $A$ は右 Noetherian である。
左の場合は反対代数に適用する。

一般の $M$ については、補題~\ref{lem:patch-complete-lifting} の加群版で
$M^{\mathrm{deg}\wedge}$ の次数付き Noetherian 性を得て、
任意の次数付き部分加群 $N\subset M$ に上と同じ有限生成の降下を適用する。
最後の主張は補題~\ref{lem:patch-graded-to-ungraded} による。
\end{proof}

\begin{corollary}[閉集合だけを指定した形]
$Z\subset\operatorname{Spec} R$ を閉集合とし、$I$ をその被約閉部分の定義イデアルとする。
$A/IA$ が右 Noetherian であり、$\operatorname{Spec} R\setminus Z$ を覆う有限個の
主開集合上で $A$ が右 Noetherian なら、$A$ は右 Noetherian である。
\end{corollary}

\begin{proof}
有限個の主開集合を $D(g_1),\ldots,D(g_s)$ と書く。
$J=(g_1,\ldots,g_s)$ とすれば $\sqrt J=I$ である。
$R$ が Noetherian なので、ある $m$ について $I^m\subset J$ である。
$A/JA$ は中央の冪零イデアル $IA/JA$ を持ち、その商は $A/IA$ である。
各 $(I^pA+JA)/(I^{p+1}A+JA)$ は有限個の中央係数を用いて
$A/IA$ の有限直和の商となるので、右 Noetherian である。
この有限フィルトレーションにより $A/JA$ は右 Noetherian である。
定理~\ref{thm:patch-noetherian} を $J$ に適用する。
\end{proof}

\subsection{Z 代数の各 row への適用}

\begin{definition}
$C$ を、$C_{ij}=0$ for $j<i$、恒等元 $e_i\in C_{ii}$、
積 $C_{ij}\otimes_R C_{jk}\to C_{ik}$ を持つ中央 $R$ 線形 $\mathbb Z$ 代数とする。
右加群は $M=\bigoplus_jM_j$ と作用 $M_j\otimes_R C_{jk}\to M_k$ からなり、
部分加群はこの頂点分解を保つものとする。
$P_i=e_iC$、すなわち $(P_i)_j=C_{ij}$ と置く。
$C$ が右 Noetherian であるとは、すべての $P_i$ が
この右加群圏で Noetherian であることをいう。
これは有限生成右 $C$ 加群がすべて Noetherian であることと同値である。
\end{definition}

\begin{theorem}[$\mathbb Z$ 代数の形式的貼合せ]
\label{thm:patch-z-algebra}
$R$ と $I=(f_1,\ldots,f_r)$ は定理~\ref{thm:patch-noetherian} の通りとし、
$C$ を各 $C_{ij}$ が有限 $R$ 加群である中央 $R$ 線形 $\mathbb Z$ 代数とする。
$C/IC$ と各 $C_{f_\alpha}$ が右 Noetherian なら、$C$ は右 Noetherian である。
左 Noetherian 性、すなわち各 column $Ce_j$ の Noetherian 性についても同様である。

各 $M_j$ が有限 $R$ 加群である右 $C$ 加群 $M$ に対しても、
$M/IM$ と各 $M_{f_\alpha}$ の Noetherian 性から $M$ の Noetherian 性が従う。
\end{theorem}

\begin{proof}
まず $R$ が $I$ 進完備な場合を考える。
補題~\ref{lem:patch-central-variable} の $\mathbb Z$ 代数版により、各 row
\[
 \overline P_i[t_1,\ldots,t_r]
 \quad\text{は}\quad
 \overline C[t_1,\ldots,t_r]
 \text{ 上 Noetherian である},
\]
ただし $\overline C=C/IC$、$\overline P_i=e_i\overline C$ である。
全射
\[
 \overline C[t_1,\ldots,t_r]\twoheadrightarrow\operatorname{gr}_I C,
 \qquad
 \overline P_i[t_1,\ldots,t_r]\twoheadrightarrow\operatorname{gr}_I P_i
\]
より、$\operatorname{gr}_I P_i$ は右 Noetherian 加群である。

$N\subset P_i$ を右部分加群とし、
$\operatorname{gr}_F N$ の有限生成元を $I$ 進次数と頂点について同次に選ぶ。
それらを
\[
 u_\ell\in N_{j_\ell}\cap I^{a_\ell}C_{i,j_\ell}
\]
に持ち上げる。固定した $x\in N_j$ の $I$ 進初項を逐次消す係数は
\[
 c_{\ell,p}\in I^{p-a_\ell}C_{j_\ell,j}
 \quad(p\geq a_\ell)
\]
に取れる。各 $C_{j_\ell,j}$ は有限 $R$ 加群なので完備であり、
係数の級数が収束する。$C_{ij}$ の分離性から
$x=\sum_\ell u_\ell c_\ell$ を得る。
よって $P_i$ のすべての右部分加群が有限生成である。
ここで $i$ ごとに有限生成元を取っており、$i$ に関する一様な生成次数の上界は不要である。

一般の $R$ では、各 $C_{ij}$ を $\widehat R_I$ に基底変換して
上の完備基底上の結論を得る。
次に \eqref{eq:patch-faithful-cover} を使い、各 $N\subset P_i$ の有限生成を
忠実平坦降下する。テンソルの生成元を展開する際に各元を頂点ごとに分ければ、
有限個の頂点同次生成元を降下できる。
したがってすべての row は Noetherian である。

一般の $M$ も、$\overline M[t_1,\ldots,t_r]\twoheadrightarrow\operatorname{gr}_I M$ を用いて同様に証明できる。
左の場合は各 column に同じ議論を適用すればよい。
\end{proof}

\begin{remark}[普通の環の Noetherian 性との区別]
無限個の直交冪等元を持つ $\bigoplus_{i,j}C_{ij}$ 全体を
単位的な Noetherian 環と呼ぶ主張ではない。
必要なのは各 row、各 column、およびその有限直和の商の Noetherian 性である。
正規族による商 $C/(g)$ を用いる場合も、閉部分上・各主開集合上で得る入力は
この row ごとの Noetherian 性でよい。
正規族の交換関係と正則性から $C/(g)$ の結果を $C$ へ持ち上げるための補題を、
次の定理~\ref{thm:patch-normal-family} で与える。
\end{remark}

\subsection{正則正規線族からの row Noetherian 性の持上げ}

\begin{theorem}[正則正規族による $\mathbb Z$ 代数の拡大]
\label{thm:patch-normal-family}
$R$ を可換環とし、$C$ を正の向きの中央 $R$ 線形 $\mathbb Z$ 代数とし、$C_{ii}=R$ とする。
整数 $\delta>0$ と可逆 $R$ 加群
$G_i\subset C_{i,i+\delta}$ $(i\in\mathbb Z)$ が与えられ、ある両側イデアル $K\subset C$ に対して
乗法が全ての $i,j$ で同型
\begin{equation}\label{eq:patch-normal-family-input}
 G_i\otimes_R C_{i+\delta,j}
       \xrightarrow{\sim}K_{ij}
       \xleftarrow{\sim}C_{i,j-\delta}\otimes_R G_{j-\delta}
\end{equation}
を与えると仮定する。
$\overline C=C/K$ が右 Noetherian なら $C$ も右 Noetherian である。
左についても同様である。
特に有限生成右加群、有限生成左加群はそれぞれ Noetherian になる。

この定理では $G_i$ の周期性、無限個の $G_i$ を同時に自明化する基底の開被覆、
各 $C_{ij}$ の $R$ 上の有限性を仮定しない。
\end{theorem}

\begin{proof}
\emph{(a) 頂点の線束 twist による正規族の自明化.}
各 $\delta$ を法とする剰余類で一つの $L_i$ を $R$ と定め、
全ての整数 $i$ について再帰的に
\begin{equation}\label{eq:patch-vertex-lines}
 L_{i+\delta}=G_i^{-1}\otimes_R L_i
\end{equation}
と定める。各固定した $i$ の $L_i$ は有限個の可逆加群の tensor 積なので、
可逆 $R$ 加群として定義される。
次の $\mathbb Z$ 代数を取る：
\[
 C'_{ij}=L_i^{-1}\otimes_R C_{ij}\otimes_R L_j.
\]
右加群の圏の完全同値は、各頂点で
$M_j\mapsto M_j\otimes_R L_j$ とすることで得られる。
この同値は $P_i$ を $L_i\otimes_R P'_i$ へ送るため、
各 row の Noetherian 性を保つ。
また
\[
 G'_i=L_i^{-1}\otimes_R G_i\otimes_R L_{i+\delta}\simeq R.
\]
これは基底環を変更する操作ではなく、頂点ごとの線束 twist である。
無限個の線束の同時 Zariski 自明化を仮定していない。

従って $C'$ を $C$ と書き直し、全ての $G_i$ に生成元
$g_i\in C_{i,i+\delta}$ がある場合を証明すればよい。
\eqref{eq:patch-normal-family-input} により、左からの $g_i$ 倍と
右からの $g_j$ 倍はそれぞれ単射で、同じ正規イデアルを与える。

\emph{(b) 添字を $\delta$ 移す自己同型.}
$a\in C_{ij}$ に対して、正規性と左からの正則性により一意な
$\sigma(a)\in C_{i+\delta,j+\delta}$ があり、
\begin{equation}\label{eq:patch-normal-shift}
 g_i\sigma(a)=a g_j
\end{equation}
を満たす。
左右の同型 \eqref{eq:patch-normal-family-input} から各 $\sigma:C_{ij}\to C_{i+\delta,j+\delta}$ は同型である。
左から $g_i$ を消去すると
$\sigma(ab)=\sigma(a)\sigma(b)$ が従い、$\sigma(e_i)=e_{i+\delta}$ である。
従って $\sigma$ は頂点を $\delta$ だけ移す $R$ 線形自己同型である。
さらに $a=g_i$ を \eqref{eq:patch-normal-shift} に代入して左から消去すれば
$\sigma(g_i)=g_{i+\delta}$ となる。
従って $\sigma$ は $K$ を保ち、$\overline C$ の自己同型も誘導する。

\emph{(c) 正の skew polynomial $\mathbb Z$ 代数の row Hilbert 基底論法.}
$D=\overline C[t;\sigma]$ を
\[
 D_{ij}=\bigoplus_{p\geq0}\overline C_{i,j-p\delta}\,t^p
\]
と定める。ここで $a t^p$ は、係数 $a\in\overline C_{i,j-p\delta}$ に
頂点 $j-p\delta$ から $j$ への形式的な $p$ 個の $t$ を付した記号である。
積は
\[
 (a t^p)(b t^q)=a\sigma^{-p}(b)t^{p+q}
\]
で定める。$\overline C$ が正の向きなので、固定した $i,j$ について
和に現れる $p$ は有限個である。

$N\subset e_iD$ を右部分加群とする。
各 $n\geq0$ に対して $L_n\subset e_i\overline C$ を、
頂点 $\ell$ の成分が
\[
 (L_n)_\ell=
 \left\{a\in\overline C_{i\ell}:
 \begin{array}{l}
 N_{\ell+n\delta}\text{ のある元が }t\text{ 次数 }n\text{ 以下であり、}\\
 \text{その }t^n\text{ の係数が }a\text{ である}
 \end{array}\right\}
\]
となる部分加群として定める。
これは実際に $\overline C$ 部分加群である。
$a\in(L_n)_\ell$ と $c\in\overline C_{\ell m}$ に対し、
$a$ の持上げに右から $\sigma^n(c)$ を掛ければ、最高係数は $ac$ となるためである。
また右から $t$ を掛けることで $L_n\subset L_{n+1}$ を得る。
$e_i\overline C$ は Noetherian なので、この昇鎖はある $n_0$ から安定する。

$0\leq n\leq n_0$ の各 $L_n$ の有限個の頂点同次生成元を取り、
$N$ の $t$ 次数 $n$ 以下の元へ持ち上げる。
これら有限個の元が $N$ を生成する。
実際、$x\in N_j$ の最高 $t$ 次数を $d$ とする。
$d\leq n_0$ なら $L_d$ の生成元を用い、
$d>n_0$ なら $L_{n_0}$ の持上げを右から $t^{d-n_0}$ 倍して用いる。
最高係数の表示に現れる $c\in\overline C_{\ell,j-d\delta}$ に対し
右から掛ける係数を $\sigma^d(c)$ とすれば、最高係数が正しく消える。
差の $t$ 次数が下がるので、有限回の次数帰納法が完了する。
よって $D$ は各 row について右 Noetherian である。

\emph{(d) $K$ 進次数付き代数.}
$g_i^{(p)}=g_i g_{i+\delta}\cdots g_{i+(p-1)\delta}$ と置く。
正規性から
\[
 K^p_{ij}=g_i^{(p)}C_{i+p\delta,j}
          =C_{i,j-p\delta}g_{j-p\delta}^{(p)}.
\]
両側の正則性により、固定した $p$ の右辺を $K^{p+1}_{ij}$ で割ることは
係数を $K_{i,j-p\delta}$ で割ることと一致する。
従って
\[
 \overline C_{i,j-p\delta}\xrightarrow{\sim}
 K^p_{ij}/K^{p+1}_{ij},
 \qquad
 \overline a\longmapsto
 \overline{a g_{j-p\delta}^{(p)}}
\]
を得る。式 \eqref{eq:patch-normal-shift} が積の整合性を与えるので、
\[
 \operatorname{gr}_K C\simeq\overline C[t;\sigma]=D
\]
である。

\emph{(e) 元の row への有限回の持上げ.}
$N\subset P_i$ に $F^pN=N\cap e_iK^p$ を入れる。
$\operatorname{gr}_F N\subset e_i\operatorname{gr}_K C$ は右有限生成である。
その有限個の $K$ 進次数と頂点についての同次生成元を
$u_1,\ldots,u_s\in N$ に持ち上げる。
固定した $x\in N_j$ の初項をこれらの右倍で逐次消去する。
正の向きと $\delta>0$ から
\[
 K^p_{ij}=0\qquad(p\delta>j-i)
\]
なので、この操作は有限回で終了し、$x\in\sum_\ell u_\ell C$ となる。
従って $P_i$ の全ての部分加群が有限生成であり、$C$ は右 Noetherian である。

左の主張は、添字を反転した反対 $\mathbb Z$ 代数
$C^{\mathrm{op,ref}}_{ij}=C_{-j,-i}$ に右の結果を適用して得る。
この変換は正の向きと、左右両方の
\eqref{eq:patch-normal-family-input} を保つ。
\end{proof}

\begin{remark}[偶数環による証明との関係]
二周期の族では、偶数部分から得る通常の次数付き環で
正則正規な次数二の線束を扱う証明も可能である。
定理~\ref{thm:patch-normal-family} はその周期性を仮定しないので、
根を選んでいない geometric cubic foundation の族にも適用できる。
必要な幾何的入力は、相対曲線商 $\overline C$ の row Noetherian 性と、
正則正規族の普遍的同型 \eqref{eq:patch-normal-family-input} である。
この二つをファイバーだけでなく実際の族として証明する必要は残る。
\end{remark}

\subsection{適用時に確認すべき入力}

閉部分 $Z\subset U=\operatorname{Spec} R$ と開補集合 $U\setminus Z$ に分ける適用では、
必要な入力は次の有限個の主張である。
\begin{enumerate}
 \item $R$ は Noetherian であり、すべての構造成分 $A_d$ または $C_{ij}$ は有限 $R$ 加群である。
 \item 閉部分を定義する実際の商環 $R/I$ 上で、$A/IA$ または $C/IC$ が Noetherian である。
 \item 開補集合を覆う有限個の主開集合 $D(f_\alpha)$ 上で、代数全体または各 row が Noetherian である。
\end{enumerate}
この三条件から結論は上の定理で従う。
Noetherian 性を各点の局所環で証明した後、任意のイデアルの生成元が
自動的に一点の近傍で生成を続ける、と推論する必要はない。
形式的貼合せは、閉部分に沿う完備化という一つの環上で全ての次数を扱い、
有限個の開局所化と合わせた忠実平坦被覆で有限生成を降下する。

